\input{preamble}
\begin{multicols}{2}[\section{Other chapters}]
\noindent
Preliminaries
\begin{enumerate}
\item \hyperref[sets-section-phantom]{Sets}
\end{enumerate}
Schemes
\begin{enumerate}
\item \hyperref[spaces-section-phantom]{Spaces}
\end{enumerate}
\end{multicols}


\title{Sets}

\begin{document}

\maketitle

\phantomsection
\label{section-phantom}

\tableofcontents

\section{Introduction}
\label{section-introduction}

\noindent
This chapter is a test file for the adapter of math2. It is in the style of
the Stacks project, and all the words in it are new words.

\section{Finite sets}
\label{section-finite}

\noindent
A set is finite when a bijection to a set of the form $\{1, \ldots, n\}$
exists.

\begin{definition}
\label{definition-finite}
A set $S$ is {\it finite} if there is a natural number $n$ and a bijection
between $S$ and the set of the natural numbers from one to $n$.
\end{definition}

\begin{lemma}[Pigeonhole]
\label{lemma-pigeonhole}
Let $S$ and $T$ be finite sets with more elements in $S$ than in $T$.
Then each map from $S$ to $T$ sends two different elements to one element.
\end{lemma}

\begin{proof}
Use induction on the number of elements of $T$. When $T$ is empty, the
set $S$ has an element and there is no map at all. In the other case we
remove one element from $T$ and look at the elements of $S$ that go to it.
\begin{proof}
An inner proof stays in the proof of the lemma.
\end{proof}
\end{proof}

\noindent
The count in the next equation is used below.
\begin{equation}
\label{equation-count}
|S \times T| = |S| \cdot |T|
\end{equation}

\subsection{Unions}
\label{subsection-unions}

\noindent
The union of two finite sets is finite.

\begin{remark}
\label{lemma-wrong-prefix}
The label of this remark starts with lemma, and the kind is still remark.
\end{remark}

\begin{lemma}
\label{lemma-items}
The following are true for finite sets $S$ and $T$.
\begin{enumerate}
\item
\label{item-union}
The union of $S$ and $T$ is finite.
\item
\label{item-product}
The product of $S$ and $T$ is finite.
\end{enumerate}
\end{lemma}

% \begin{lemma}
% \label{lemma-comment}
% A lemma in a comment is not a lemma.
% \end{lemma}

\begin{multicols}{2}[\section{Other chapters}]
\noindent
Preliminaries
\begin{enumerate}
\item \hyperref[sets-section-phantom]{Sets}
\end{enumerate}
Schemes
\begin{enumerate}
\item \hyperref[spaces-section-phantom]{Spaces}
\end{enumerate}
\end{multicols}


\end{document}
